\documentclass[11pt,a4paper]{article}

\usepackage[utf8]{inputenc}
\usepackage[T1]{fontenc}
\usepackage[french]{babel}
\usepackage[margin=2.5cm]{geometry}
\usepackage{amsmath,amssymb}
\usepackage{xcolor}
\usepackage{tcolorbox}
\tcbuselibrary{theorems,breakable}
\usepackage{enumitem}
\usepackage{fancyhdr}
\usepackage{booktabs}

% En-tête et pied de page
\pagestyle{fancy}
\fancyhf{}
\fancyhead[L]{\small TD S01 -- Qu'est-ce que l'IA ?}
\fancyhead[R]{\small Introduction à l'IA -- M1 IABD}
\fancyfoot[C]{\thepage}
\renewcommand{\headrulewidth}{0.4pt}

% Boîte pour les exercices
\newtcolorbox{exercicebox}[1][]{
    colback=blue!5,
    colframe=blue!50,
    fonttitle=\bfseries,
    title=#1,
    breakable
}

\title{\textbf{TD Séance 1 : Qu'est-ce que l'Intelligence Artificielle ?}\\
\large Introduction à l'Intelligence Artificielle --- Énoncés}
\author{Dr. El Hadji Bassirou TOURÉ\\
Département de Mathématiques et Informatique\\
Faculté des Sciences et Techniques\\
Université Cheikh Anta Diop de Dakar}
\date{M1 IABD --- Semestre 1}

\begin{document}

\maketitle

\vspace{0.5cm}

\begin{center}
\fbox{\parbox{0.9\textwidth}{
\textbf{Objectifs du TD :}
\begin{itemize}[nosep]
    \item Maîtriser la distinction IA / ML / DL
    \item Savoir appliquer le cadre PEAS à un problème
    \item Caractériser les environnements selon leurs propriétés
    \item Identifier le type d'agent approprié à une situation
    \item Réfléchir aux limites des définitions classiques de l'IA
\end{itemize}
}}
\end{center}

\vspace{1cm}

%=============================================================================
\section*{Exercice 1 : Classification IA / ML / DL \hfill \textit{(20 min)}}
%=============================================================================

\begin{exercicebox}[Énoncé]
Pour chacun des systèmes suivants, indiquez s'il relève de :
\begin{itemize}
    \item[(a)] IA mais pas ML (IA symbolique, algorithmique)
    \item[(b)] ML mais pas DL
    \item[(c)] Deep Learning
\end{itemize}

Justifiez brièvement chaque réponse.

\begin{enumerate}
    \item Un algorithme de navigation GPS qui calcule le plus court chemin
    \item Un système de recommandation Netflix utilisant un réseau de neurones à 50 couches
    \item Un filtre anti-spam utilisant Naive Bayes
    \item Wolfram Alpha résolvant des intégrales symboliquement
    \item AlphaGo battant le champion du monde de Go
    \item Un arbre de décision pour prédire si un client va résilier son abonnement
    \item Un système expert médical basé sur des règles if-then
    \item GPT-4 générant du texte
    \item Un algorithme k-means pour segmenter des clients
    \item Un solveur SAT vérifiant la satisfiabilité d'une formule logique
\end{enumerate}
\end{exercicebox}

\vspace{1cm}

%=============================================================================
\section*{Exercice 2 : Analyse PEAS \hfill \textit{(30 min)}}
%=============================================================================

\begin{exercicebox}[Énoncé]
Définissez le cadre PEAS (Performance, Environment, Actuators, Sensors) pour les agents suivants. Pour chaque composante, donnez au moins 3-4 éléments pertinents.

\begin{enumerate}
    \item Un robot aspirateur autonome (type Roomba)
    \item Un système de trading haute fréquence
    \item Un assistant vocal (type Alexa/Siri)
\end{enumerate}

Pour chaque agent, précisez également :
\begin{itemize}
    \item Quel compromis existe entre les différents critères de performance ?
    \item Quelle est la principale difficulté liée à l'environnement ?
\end{itemize}
\end{exercicebox}

\vspace{1cm}

%=============================================================================
\section*{Exercice 3 : Types d'environnements \hfill \textit{(25 min)}}
%=============================================================================

\begin{exercicebox}[Énoncé]
Pour chaque environnement ci-dessous, caractérisez-le selon les 6 dimensions vues en cours :
\begin{itemize}
    \item Observable : Total / Partiel
    \item Déterministe : Oui / Non
    \item Épisodique : Oui / Non
    \item Statique : Oui / Non
    \item Discret : Oui / Non
    \item Nombre d'agents : Un / Plusieurs
\end{itemize}

Complétez le tableau suivant et justifiez chaque réponse :

\begin{center}
\begin{tabular}{|l|c|c|c|c|c|c|}
\hline
\textbf{Environnement} & \textbf{Obs.} & \textbf{Dét.} & \textbf{Épis.} & \textbf{Stat.} & \textbf{Disc.} & \textbf{Agents} \\
\hline
Le jeu du morpion & & & & & & \\
\hline
La bourse (pour un trader) & & & & & & \\
\hline
Un examen QCM en ligne & & & & & & \\
\hline
Une partie de poker Texas Hold'em & & & & & & \\
\hline
Un match de football (robot joueur) & & & & & & \\
\hline
\end{tabular}
\end{center}

Questions complémentaires :
\begin{enumerate}
    \item Quel environnement est le plus facile à gérer pour un agent ? Pourquoi ?
    \item Quel environnement est le plus difficile ? Pourquoi ?
    \item Pour la bourse, expliquez pourquoi elle est « partiellement observable » même si on a accès aux prix en temps réel.
\end{enumerate}
\end{exercicebox}

\vspace{1cm}

%=============================================================================
\section*{Exercice 4 : Type d'agent approprié \hfill \textit{(25 min)}}
%=============================================================================

\begin{exercicebox}[Énoncé]
Pour chaque situation, déterminez le type d'agent \textbf{le plus simple} qui pourrait la gérer efficacement. Choisissez parmi :
\begin{itemize}
    \item Agent réflexe simple
    \item Agent basé sur modèle (réflexe avec état)
    \item Agent basé sur buts
    \item Agent basé sur utilité
    \item Agent apprenant
\end{itemize}

Justifiez votre choix en expliquant :
\begin{itemize}
    \item Pourquoi un agent plus simple ne suffirait pas
    \item Pourquoi un agent plus complexe serait inutilement sophistiqué
\end{itemize}

\textbf{Situations :}
\begin{enumerate}
    \item Un distributeur automatique de boissons
    \item Un thermostat intelligent qui apprend les préférences de l'utilisateur
    \item Un système de recommandation de films
    \item Un planificateur de trajet GPS
    \item Une porte automatique dans un supermarché
    \item Un robot qui joue aux échecs contre des adversaires de niveaux variés
    \item Un système de contrôle de feux de circulation adaptatif
\end{enumerate}
\end{exercicebox}

\vspace{1cm}

%=============================================================================
\section*{Exercice 5 : Réflexion -- Limites du test de Turing \hfill \textit{(20 min)}}
%=============================================================================

\begin{exercicebox}[Énoncé]
Un système réussit le test de Turing : lors d'une conversation textuelle de 5 minutes, 30\% des juges humains le confondent avec un humain (seuil traditionnellement accepté).

\textbf{Questions :}
\begin{enumerate}
    \item \textbf{Ce système est-il nécessairement intelligent ?} 
    
    Argumentez en considérant :
    \begin{itemize}
        \item L'argument de la « Chambre chinoise » de John Searle
        \item Le cas du programme ELIZA (1966)
        \item La distinction entre « simuler » et « comprendre »
    \end{itemize}
    
    \item \textbf{Un système peut-il être intelligent sans réussir le test de Turing ?}
    
    Donnez au moins trois exemples de systèmes d'IA reconnus comme performants qui échoueraient probablement au test de Turing.
    
    \item \textbf{Quelles capacités le test de Turing ne mesure-t-il pas ?}
    
    Listez au moins 5 capacités considérées comme « intelligentes » mais non évaluées par ce test.
    
    \item \textbf{Question ouverte :} Proposez un test alternatif qui vous semblerait plus approprié pour évaluer l'intelligence d'un système artificiel. Justifiez votre proposition.
\end{enumerate}
\end{exercicebox}

\vspace{1cm}

%=============================================================================
\section*{Exercice 6 (Bonus) : Conception d'un agent \hfill \textit{(30 min)}}
%=============================================================================

\begin{exercicebox}[Énoncé]
Vous devez concevoir un agent pour un \textbf{robot livreur de colis} dans un campus universitaire.

\textbf{Partie A -- Analyse PEAS}
\begin{enumerate}
    \item Définissez complètement le cadre PEAS pour cet agent.
    \item Identifiez les conflits potentiels entre les critères de performance.
\end{enumerate}

\textbf{Partie B -- Caractérisation de l'environnement}
\begin{enumerate}
    \item Caractérisez l'environnement selon les 6 dimensions.
    \item Quelles sont les principales sources d'incertitude ?
    \item Comment l'environnement varie-t-il entre le jour et la nuit ?
\end{enumerate}

\textbf{Partie C -- Architecture de l'agent}
\begin{enumerate}
    \item Quel type d'agent recommanderiez-vous ? Justifiez.
    \item Décrivez les principales composantes de cet agent :
    \begin{itemize}
        \item État interne (si applicable)
        \item Règles ou politique de décision
        \item Mécanisme d'apprentissage (si applicable)
    \end{itemize}
    \item Donnez un exemple de situation où l'agent devrait :
    \begin{itemize}
        \item Utiliser son modèle de l'environnement
        \item Réviser ses croyances
        \item Faire un compromis entre ses objectifs
    \end{itemize}
\end{enumerate}
\end{exercicebox}

\end{document}
